\scanpage{604}\mspage{13}
% archive sheet number: 603

\[
\underbrace{\underline S'=S'_{\unsure{a|\cdot}}}_{\Vert\,\text{\unsure{def|déf}}}\subset M
\]
\[
\begin{aligned}
 f^{-1}(\{\emptyset\})={}&\text{ss-cat. \deleted{\unsure{str. pleine de|fermée de}} }M\\
 &\text{formée des }x\in\mathrm{Ob}\,M\text{ tels}\\
 &\text{que } a(x)=\emptyset .
\end{aligned}
\]

Comme $\underline S'$, \deleted{\unsure{étant}} fermée dans $M$, est
stable par $\unsure{\varinjlim|\underline{\lim}}$ non vides et en
particulier $L_{\pi}$-stable pour $\pi$-\unsure{gros|gd}, il y a plus de raisons de
s'attendre à ce que $\underline S'$ soit accessible dans $M$ que pour
\inserted{son complémentaire} $\underline S$, \deleted{\unsure{qui}}
dont il n'est pas clair qu'il ait cette propriété de stabilité, mais
seulement stabilité par $\varprojlim$ non vides, et en particulier par
facteurs directs.

En fait, nous allons voir qu'il peut arriver que $\underline S$ et
$\underline S'$ soient tous deux non accessibles dans $M$ :

\uline{Proposition} \uline{4} \unsure{Soient|Soit} $\mathcal U$ un univers, contenant
un cardinal infini $\pi_0$, tel que pour tt $E\in\mathcal U$,
\deleted{\unsure{tt}} \deleted{\unsure{tout}} $\pi_0$-ultrafiltre sur
$E$ soit trivial. Soit alors $k\in\mathcal U$ un corps de cardinal
\inserted{$\pi'_0$}$\geq\pi_0$, et soit $M$ la catégorie
